\section{Ingredients, missing chemistry, and data resources}
The next simulator components should be selected by whether competition
measurements can constrain them. Our first priorities are free magnesium,
explicit nucleotide salvage and energy chemistry, and target-specific folding.
The linked draft catalog already contains magnesium/potassium glutamate,
HEPES, phosphate buffers, NTPs, NMPs, nucleosides, ribose, glucose, PEP, pyruvate,
acetyl phosphate, maltose, maltodextrin, NAD, and CoA. It also lists DsbC,
reduced/oxidized glutathione, pyruvate oxidase, and catalase\cite{competition_rules}.
These additions remain distinct from unknown endogenous pools. The draft's
oxidized-glutathione description repeats the reduced abbreviation; reconcile
chemical identity with the stated supplier product before interpreting redox.

Model magnesium binding to nucleotides and phosphate before treating total
magnesium as free magnesium. Track AMP, ADP, GDP, PPi, and Pi rather than hiding
all energy chemistry in an ATP-regeneration term. Represent each energy substrate
through the enzymes and cofactors that actually convert it; the abstract Fuel
variable cannot establish that glucose and PEP are interchangeable. Add oxygen
transfer and peroxide removal when considering the pyruvate-oxidase/catalase
module. For enzyme targets, distinguish synthesized, soluble, folded, and active
protein. DsbC and glutathione are already catalog options; chaperone systems such
as DnaK/DnaJ/GrpE and GroEL/GroES are possible additional requests, motivated
by published CFPS folding work\cite{hunt}. Test functional yield and assay-matrix
effects rather than assume that brighter GFP establishes better enzyme production.

Other possible requests are nucleoside-diphosphate kinase, adenylate kinase, and
inorganic pyrophosphatase, to test nucleotide replenishment and pyrophosphate
removal. These are proposed modules, not validated improvements. Sourcing,
stock activity, storage formulation, and competition acceptance need verification.
The draft allows requests for additional commercially available reagents but
leaves the team allotment as ``10 or 20''; it is not a confirmed quota
\cite{competition_rules}. The machine-readable
\texttt{examples/research\_priorities.json} separates current draft examples
from proposed requests and lists missing mechanisms. No new reagent is silently
assigned a kinetic effect in the present simulator.

\subsection{Databases for an ingredient and mechanism registry}
\begin{center}
\small
\begin{tabular}{p{0.24\linewidth}p{0.65\linewidth}}
Resource & Use in this project\\\hline
ChEBI and PubChem & Normalize identities, structures, synonyms, and salt forms;
retain stock-specific composition and units\cite{chebi,pubchem}.\\
Rhea & Retrieve curated biochemical reaction participants and stoichiometry
using ChEBI identifiers\cite{rhea}.\\
SABIO-RK and BRENDA & Find kinetic laws, constants, enzyme context, and measurement
conditions; use matched values as uncertain priors\cite{sabio,brenda}.\\
eQuilibrator and NIST & Check biochemical thermodynamics and dependence on pH,
ionic conditions, and concentrations\cite{equilibrator,nist}.\\
BioNumbers & Find cited biological scales and parameter ranges, retaining entry
identifiers and underlying references\cite{bionumbers}.\\
BioModels and author repositories & Retrieve executable models and reproduce
original observations before importing mechanisms\cite{biomodels,horvath_code}.\\
\end{tabular}
\end{center}

These resources describe chemistry and mechanisms, not competition efficacy.
The competition catalog and measurements remain the authority for feasible
formulations and observed performance. For every datum, retain source/version,
organism and enzyme identity, substrate, pH, temperature, ionic conditions,
units, uncertainty, and access/reuse terms. A purified-enzyme turnover number
is not a lysate $V_{\max}$ without an active enzyme concentration. Distinguish
measured constants from estimates and inferred values. Represent salts,
counterions, hydrates, enzyme activity units, mixed amino-acid stocks, and polymer
distributions explicitly; molecular-weight lookup alone cannot convert every
stock into one meaningful molarity.

\subsection{What the choice of T7 tells us}
The draft specifies T7-driven supplied circular templates and an internally
produced BL21 Star (DE3) lysate, with IPTG used during fermentation to induce T7
RNA polymerase\cite{competition_rules}. Thus T7 polymerase is expected to be
provided by the extract; its active concentration in the final reaction is still
unknown. The supplied promoter is not an adjustable competition variable, and
an additional polymerase supplement is not automatically needed.

T7 offers promoter-specific transcription with a dedicated phage polymerase,
making transcription capacity a useful module to distinguish from native
translation capacity\cite{yang}. A strong transcription system does not guarantee
maximum protein output: mRNA production consumes NTPs and can exceed translation
capacity. Magnesium and nucleotide conditions can affect both processes. The
computed transcription/translation grid probes this resource-allocation hypothesis
in our reduced model; it cannot measure actual T7 activity. Separating polymerase
and ribosome resources is a future improvement over the shared Machinery state.
Calibrate transcript production and decay where available, then protein synthesis
and maturation, instead of absorbing every discrepancy into one transcription rate.

The promoter motif supports an assumed transcription start, not a measured
termination boundary on the circular plasmid. Readthrough changes NTP demand
while CDS amino-acid demand stays the same; our boundary scenario makes this
uncertainty visible. Borkowski et al.'s study uses a constitutive-promoter reporter
\cite{borkowski}, so its fitted transcription behavior cannot be imported unchanged
into our T7 system. Recipe-design lessons transfer as hypotheses; promoter-specific
rates need their own evidence.

\section{Argonne, scientific data infrastructure, and the October hackathon}
The researcher is Ben Blaiszik, affiliated with Globus, Data Science and Learning
at the University of Chicago and Argonne National Laboratory\cite{blaiszik}.
His data-ecosystem work links the Materials Data Facility and DLHub; Globus Labs'
Foundry adds standardized dataset metadata and structures for discovering and
connecting data with models\cite{foundry}. Here ``Foundry'' denotes data
infrastructure and is distinct from Illinois' iBioFAB wet-laboratory biofoundry.
The applicable lesson is to make measurements, processing steps, models, and
benchmarks independently discoverable and reproducible. A materials database
is not a ready-made source of validated CFPS recipes.

The shared event is now the \emph{Open Scientific Intelligence Hackathon},
formerly the LLM Hackathon for Materials Science and Chemistry. The official
site announces October 21--22, 2026, as a free hybrid international event covering
agents, datasets, benchmarks, models, and scientific software\cite{osi}.
Its scope has expanded to physical sciences and mathematics, so a CFPS project
should explain its biochemical mechanism and data/benchmark contribution;
we have not assumed acceptance or registered a project.

Our proposed contribution is an evidence-grounded CFPS recipe and calibration
assistant. It would extract literature recipes and kinetic evidence into a
structured registry with source passages, normalize chemical identities and
units, validate feasibility, and call versioned simulation/calibration tools.
Benchmark extraction on manually checked documents, conversion on known stocks,
and predictions on withheld experiments. Measure unsupported ingredient claims
and incorrect concentration conversions alongside scientific utility. An LLM
can suggest candidates for review, while numerical tools enforce actual stock
and dispensing constraints. Present an end-to-end example that retains evidence
and fails explicitly when facts are missing. This is a proposed hackathon project;
no agent-based extraction, remote data publication, or experimental scheduling
is implemented by the present suite.
